% EPPA-III, Lemma \label{lemma:gtg} (current Section 5.1).
% Source base: main-new(20261019-085906).tex.
% This is a separate proposed replacement, NOT an edit of the frozen manuscript.
% No extra theorem is cited; the Mallows--Sloane theorem may still be
% kept earlier as background/motivation.  The proof covers empty
% domains, all fixed points, even and odd cycles, and open paths.
%
% Proposed replacement for the proof of \label{lemma:gtg} only:

\begin{proof}
Choose a graph $F$ in the switching class and write
$e(x,y)\in\mathbb F_2$ for its adjacency indicator.
Put $D=\dom{f}$ and, for $x,y\in D$, let
$\delta(x,y)=e(x,y)+e(fx,fy)$.
Since $f$ preserves the two-graph, we have
$$
\delta(x,y)+\delta(y,z)+\delta(z,x)=0
$$
for every $x,y,z\in D$; the equality also holds when two
vertices coincide. If $D$ is empty there is nothing to prove.
Otherwise choose $r\in D$ and put $h(x)=\delta(r,x)$.
The displayed equality gives
$$
\delta(x,y)=h(x)+h(y)\qquad(x,y\in D).
$$

We will choose a function $s:X\to\mathbb F_2$ and a constant
$c\in\mathbb F_2$ such that
$$
s(x)+s(fx)=h(x)+c\qquad(x\in D).
$$
Indeed, switching $F$ according to $s$ changes its adjacency
indicator to $e_s(x,y)=e(x,y)+s(x)+s(y)$, and then, for
$x,y\in D$,
$$
e_s(x,y)+e_s(fx,fy)
=\delta(x,y)+h(x)+h(y)+c+c=0.
$$
Thus $f$ will be a partial automorphism of the switched graph.

The directed graph of the partial bijection $f$ is a disjoint
union of directed paths and cycles. On a cycle $C$ the displayed
equations for $s$ are consistent precisely when
$$
\sum_{x\in C}h(x)=|C|c
$$
in $\mathbb F_2$. We claim that one choice of $c$ satisfies
this condition for every cycle.

First, let $C=(x_0,\ldots,x_{2m-1})$ be an even cycle, with
$fx_i=x_{i+1}$ and indices read modulo $2m$. Summing
$\delta(x_i,x_{i+m})=e(x_i,x_{i+m})
+e(x_{i+1},x_{i+m+1})$ over $0\le i<m$
telescopes to zero: the final unordered pair is the first.
On the other hand, these $m$ pairs partition $C$, so
$$
\sum_{x\in C}h(x)
=\sum_{i=0}^{m-1}\delta(x_i,x_{i+m})=0.
$$

Next, let $C,C'$ be two odd cycles of lengths $a,b$ and choose
$x\in C$, $y\in C'$. Put $L=\operatorname{lcm}(a,b)$.
Summing $\delta(f^ix,f^iy)
=e(f^ix,f^iy)+e(f^{i+1}x,f^{i+1}y)$
over $0\le i<L$ again telescopes to zero. Hence
$$
0=\frac La\sum_{z\in C}h(z)
 +\frac Lb\sum_{z\in C'}h(z).
$$
Both coefficients are odd, so the two sums are equal.
If there is an odd cycle, choose $c$ to be this common sum;
otherwise put $c=0$. The equations are now consistent on
all cycles.

Finally choose $s$ at one point of each cycle and propagate
it using $s(fx)=s(x)+h(x)+c$; the cycle condition ensures
consistency. On each directed path choose $s$ at its initial
point freely and propagate by the same equation. Choose
$s$ arbitrarily on vertices outside $\dom{f}\cup\ran{f}$.
This constructs the required switching and proves the lemma.
\end{proof}
